\section{Точная геометрия существенных координат Rule 110}\label{sec:geometry}
\subsection{Наблюдаемая и однобитовые свидетельства}
Фиксируем одну целевую клетку в начале координат и горизонт $t$. Входная координата существенна, если найдётся начальная конфигурация, в которой изменение одного этого бита меняет ответ. Фон может зависеть как от координаты, так и от горизонта. Это различие кванторов отделяет предельную чувствительность от скорости конкретного глайдера. Для Rule~110 известны конструкция универсального вычисления \cite{cook2004} и классификация движущихся структур \cite{wolfram1986,martinez2006}. Связь наблюдений с предсказательными классами и ограничениями вывода обсуждается в \cite{shalizi2001,wolpert2008}.

\begin{definition}[Rule 110 и целевая наблюдаемая]
Положим $X = \{0, 1\}^{\mathbb{Z}}$ и зададим локальное правило над полем $\mathbb{F}_2$ алгебраической нормальной формой
\begin{equation}\label{eq:rule110}
f(l, c, r) = c \oplus r \oplus cr \oplus lcr.
\end{equation}
Синхронный оператор эволюции $F : X \to X$ определяется равенством $(F(u))_i = f(u_{i-1}, u_i, u_{i+1})$. Наблюдаемая в начале координат в момент времени $t \in \mathbb{N}$ есть $G_t(u) = (F^t(u))_0$.
\end{definition}

\begin{definition}[Предсказательная факторизация и существенные переменные]
Отображение наблюдения $O : X \to \mathcal{O}$ называется достаточным для $G_t$, если существует отображение $P : \mathcal{O} \to \{0, 1\}$ такое, что $G_t = P \circ O$. Для $j \in \mathbb{Z}$ определим однобитовое возмущение $(u \oplus e_j)_i = u_i \oplus \mathbf{1}_{i=j}$. Множество существенных начальных координат определяется как
\begin{equation}\label{eq:essential_set}
D_t = \{ j \in \mathbb{Z} : \exists u \in X,\; G_t(u) \ne G_t(u \oplus e_j) \}.
\end{equation}
\end{definition}

Поскольку $G_t$ зависит лишь от конечного числа координат, $D_t$ является единственным минимальным по включению множеством входных координат, достаточным для вычисления $G_t$, а любой достаточный связный интервал наблюдения обязан содержать целочисленную выпуклую оболочку множества $D_t$.

Для работы со свидетельствами на периодических фонах удобно фиксировать связь между входной координатой $j$ и физической координатой $s$ дефекта на решётке. Для периодического слова $w$, пространственной фазы $\phi$ и соответствующей бесконечной конфигурации $B_0^{w,\phi}$ положим
\begin{equation}\label{eq:diff_support}
B_t^{w,\phi} = F^t(B_0^{w,\phi}), \qquad
\widetilde B_t^{w,\phi} = F^t(B_0^{w,\phi} \oplus e_0), \qquad
S_t^{w,\phi} = \{ s \in \mathbb{Z} : B_t^{w,\phi}(s) \ne \widetilde B_t^{w,\phi}(s) \}.
\end{equation}
В силу сдвиговой эквивариантности $F \circ \sigma^d = \sigma^d \circ F$, где $(\sigma^d u)_i = u_{i-d}$, любая точка $s \in S_t^{w,\phi}$ порождает свидетельство $j = -s \in D_t$ после сдвига начальной пары на $-s$. Следовательно, для любого семейства фонов и фаз $\mathcal{W}$ выполняется включение
\begin{equation}\label{eq:coverage_bridge}
C_t(\mathcal{W}) = \bigcup_{(w,\phi) \in \mathcal{W}} \{ -s : s \in S_t^{w,\phi} \} \subseteq D_t.
\end{equation}
При этом левая граница входного окна $j = -L(t)$ соответствует правому физическому фронту дефекта $s = L(t)$, а правая граница входного окна $j = t$ — левому физическому фронту $s = -t$.

\subsection{Точное окно зависимости на бесконечной решётке}

Положим $L(t) = t - \lfloor t/3 \rfloor$ и $I_t = [-L(t), t] \cap \mathbb{Z}$.

\begin{lemma}[Трёхтактное правостороннее исключение]\label{lem:three_step}
Для любого семибитного слова $b = (b_0, \dots, b_6) \in \{0, 1\}^7$ обозначим через $Q(b)$ выход центрального узла после трёх синхронных шагов Rule 110. Тогда
\begin{equation}\label{eq:seven_bit}
Q(b_0, b_1, \dots, b_6) = Q(1 \oplus b_0, b_1, \dots, b_6).
\end{equation}
\end{lemma}

\begin{proof}
Из \eqref{eq:rule110} булева производная по левому аргументу равна
\begin{equation}\label{eq:bool_deriv}
f(1, c, r) \oplus f(0, c, r) = cr.
\end{equation}
Пусть $b_i(t)$ обозначает значение клетки $i$ в одной из рассматриваемых траекторий. Предположим, что две конфигурации строго совпадают правее узла $p$, и их крайнее правое различие максимально смещается вправо за три такта: $p \to p+1 \to p+2 \to p+3$. Первый сдвиг вправо требует, чтобы общие значения в $p+1$ и $p+2$ при $t=0$ равнялись $1$; второй сдвиг требует значений $1$ в $p+2$ и $p+3$ при $t=1$. Поскольку $f(1, 1, z) = 1 \oplus z$, это вынуждает $b_{p+3}(0) = 0$. Далее, так как $f(1, 0, z) = z$, для поддержания фронта необходимо $b_{p+4}(0) = 1$. Но тогда $b_{p+4}(1) = f(0, 1, b_{p+5}(0)) = 1$, откуда $b_{p+3}(2) = f(1, 1, 1) = 0$, что противоречит третьему максимальному сдвигу вправо. Тождество \eqref{eq:seven_bit} также исчерпывающе проверено на всех $2^7 = 128$ входах в ядре Lean 4.
\end{proof}

\begin{theorem}[Универсальное включение множества зависимости]\label{thm:upper_bound}
Для каждого $t \in \mathbb{N}$ справедливо включение $D_t \subseteq I_t = [-(t - \lfloor t/3 \rfloor), t] \cap \mathbb{Z}$.
\end{theorem}

\begin{proof}
Если две конфигурации совпадают строго правее $p_0$, то за один шаг согласие сохраняется строго правее $p_0 + 1$, а по Лемме~\ref{lem:three_step} за три шага — строго правее $p_0 + 2$. Представляя $t = 3q + r$ с $q = \lfloor t/3 \rfloor$ и $r \in \{0, 1, 2\}$, получаем сохранение согласия правее $p_0 + 2q + r = p_0 + L(t)$. Следовательно, различие в точке $j < -L(t)$ не достигает начала координат к моменту $t$. При $j > t$ распространение блокируется стандартным световым конусом радиуса $1$.
\end{proof}

\begin{theorem}[Достижимость обеих границ]\label{thm:sharpness}
Для каждого $t \in \mathbb{N}$ обе крайние точки $-L(t)$ и $t$ принадлежат $D_t$. Минимальный достаточный связный интервал наблюдения равен $I_t$ и имеет ширину
\begin{equation}
m_{\min}(t) = 2t + 1 - \lfloor t/3 \rfloor, \qquad \lim_{t \to \infty} \frac{m_{\min}(t)}{t} = \frac{5}{3}.
\end{equation}
\end{theorem}

\begin{proof}
Для правой границы $j = t$ достаточно сравнить нулевую конфигурацию $0^{\mathbb{Z}}$ с одиночной единицей $e_t$: в силу $f(0, 0, 0) = 0$ и $f(0, 0, 1) = 1$ крайняя левая единица смещается ровно на одну клетку влево за каждый такт и достигает нуля за $t$ шагов.

Для левой границы $j = -L(t)$ рассмотрим фон периода четыре $B_0(i) = (1, 0, 1, 1)[i \bmod 4]$ и возмущение $\widetilde B_0 = B_0 \oplus e_2$. Прямая итерация Rule 110 на шести тактах даёт эволюцию носителя различий $S_t$, представленную в таблице~\ref{tab:period4_defect}.

\begin{table}[h]
\centering
\caption{Эволюция одиночного дефекта на фоне $(1011)$ и периодический сдвиг с шагом $(\Delta t, \Delta s) = (3, +2)$}
\label{tab:period4_defect}
\begin{tabular}{cccc}
\toprule
$t$ & Слово фона $B_t$ (период $4$) & Носитель различий $S_t$ & $\max S_t$ \\
\midrule
$0$ & $1011$ & $\{2\}$ & $2$ \\
$1$ & $1110$ & $\{1, 3\}$ & $3$ \\
$2$ & $1011$ & $\{1, 3, 4\}$ & $4$ \\
$3$ & $1110$ & $\{0, 1, 4\}$ & $4$ \\
$4$ & $1011$ & $\{-1, 0, 1, 3, 5\}$ & $5$ \\
$5$ & $1110$ & $\{3, 5, 6\}$ & $6$ \\
$6$ & $1011$ & $\{2, 3, 6\}$ & $6$ \\
\bottomrule
\end{tabular}
\end{table}

Из локальности эволюции и совпадения срезов при $t=3$ и $t=6$ со сдвигом на $2$ клетки вправо следуют точные тождества конфигураций $B_6 = \sigma^2 B_3$ и $\widetilde B_6 = \sigma^2 \widetilde B_3$. По сдвиговой эквивариантности $(B_{3+3k+r}, \widetilde B_{3+3k+r}) = \sigma^{2k}(B_{3+r}, \widetilde B_{3+r})$ для всех $k \ge 0$ и $r \in \{0, 1, 2\}$, откуда $\max S_t = 2 + L(t)$ для всех $t \ge 0$. Сдвиг начальной пары на $-[2 + L(t)]$ переносит начальное возмущение в $-L(t)$, а различие в момент $t$ — в $0$, что доказывает $-L(t) \in D_t$.
\end{proof}

\subsection{Конечные кольца и расхождение между ёмкостью и доступом}

На конечном периодическом кольце $\mathbb{Z}_N$ ($N \ge 3$) с пространством конфигураций $X_N = \{0, 1\}^{\mathbb{Z}_N}$ и оператором $F_N : X_N \to X_N$ обозначим через $\mathcal{R}_t^{(N)} = F_N^t(X_N)$ множество достижимых состояний, а через $C_t^{(N)} = \lceil \log_2 |\mathcal{R}_t^{(N)}| \rceil$ — битовую ёмкость состояния. Введём спаренную динамику $\Phi_N(u, v) = (F_N(u), F_N(v))$, начальные пары $E_j = \{ (u, u \oplus e_j) : u \in X_N \}$ и множество различий в начале координат $W_0 = \{ (u,v) : u_0 \ne v_0 \}$. Тогда
\begin{equation}
j \in D_t^{(N)} \iff \Phi_N^t(E_j) \cap W_0 \ne \emptyset.
\end{equation}

\begin{proposition}[Конечный сертификат персистентной существенности]\label{prop:ring_cert}
Пусть $\mu_N$ — максимальная длина предпериода конечной детерминированной системы $F_N$, а $P_N$ — наименьшее общее кратное длин её циклов. Тогда $D_{t+P_N}^{(N)} = D_t^{(N)}$ при всех $t \ge \mu_N$. Следовательно, для любого $T \in \mathbb{N}$ условие $\forall t \ge T,\; D_t^{(N)} = \mathbb{Z}_N$ эквивалентно проверке $\Phi_N^t(E_j) \cap W_0 \ne \emptyset$ для всех $j \in \mathbb{Z}_N$ и всех $T \le t < \max(T, \mu_N) + P_N$.
\end{proposition}

\begin{proof}
Каждое достижимое состояние $F_N^t(u)$ при $t\ge\mu_N$ лежит на цикле, длина которого делит $P_N$. Следовательно, отображения $F_N^{t+P_N}$ и $F_N^t$ совпадают на всём $X_N$. Для любого $j$ это сохраняет совпадение или различие во всех парах $(u,u\oplus e_j)$. С момента $\max(T,\mu_N)$ достаточно проверить один период. Если $T<\mu_N$, дополнительно проверяются все времена от $T$ до $\mu_N-1$.
\end{proof}

\begin{remark}[Компрессия состояний и чувствительность к входам]
Локальная чувствительность на бесконечной решётке не гарантирует персистентной существенности на конечном кольце: для Rule 110 при $N = 5$ все $32$ состояния схлопываются в $0^5$ к моменту $t = 5$, давая $D_t^{(5)} = \emptyset$ при $t \ge 5$. Напротив, полный перебор для Rule 110 на кольце $N = 20$ при $t = 121$ обнаруживает $|\mathcal{R}_{121}^{(20)}| = 2095$ достижимых состояний ($12$ бит ёмкости хранения), но при этом $D_{121}^{(20)} = \mathbb{Z}_{20}$ — все $20$ входных бит остаются существенными. Это разделяет четыре независимых аспекта задачи предсказания: \emph{причинность} (носитель $D_t$), \emph{доступ} (индивидуальное чтение координат против предобработки), \emph{ёмкость} ($C_t^{(N)}$) и \emph{вычислительную сложность}.
\end{remark}

\begin{theorem}[Сплошной носитель существенности]\label{thm:gap-free}
Для каждого $t \in \mathbb{N}$ выполняется равенство $D_t = I_t = [-(t - \lfloor t/3 \rfloor), t] \cap \mathbb{Z}$.
\end{theorem}

Первые конечные проверки охватили $0\le t\le2000$. Более поздний эксперимент с фиксированными фазами семи слов продолжил проверку до $t=5000$. Доказательство для всех времён завершается в подразделе~\ref{sec:five-bands}. Там используется новый конечный сертификат с границей $t=2484$.

\subsection{Односторонние автономные тайлы и две бесконечные приграничные полосы}

Из Леммы~\ref{lem:three_step} композицией вытекает, что область зависимости $3m$-тактного выхода в ячейке $s$ удовлетворяет включению $\mathrm{Dep}_{3m}(s) \subseteq [s - 2m, s + 3m]$ для любых конфигураций.

\begin{lemma}[Одностороннее полулучевое замыкание]\label{lem:one_sided}
Пусть $A_t = F^t(A_0)$, $p > 0$ и $a(t + p) = a(t) + \delta$. Если на базовых классах $t_0 \le t < t_0 + p$ выполнено $A_{t+p}(s + \delta) = A_t(s)$ для всех $s \ge a(t)$, и вычисление $A_{t+p}$ на полулуче $[a(t+p), \infty)$ зависит только от значений $A_t$ на $[a(t), \infty)$, то сдвиговое равенство $A_{t+p}(s + \delta) = A_t(s)$ сохраняется для всех $t \ge t_0$ и $s \ge a(t)$. Аналогичное утверждение верно для левых полулучей $s \le a(t)$.
\end{lemma}

Применяя Лемму~\ref{lem:one_sided} к правому физическому фронту $e(t) = L(t)$ с шагом $(\Delta t, \Delta s) = (30, +20)$ (где $e(t+30) = e(t) + 20$, а область зависимости нового тайла ширины $W$ лежит в $[e(t) - W + 1, e(t) + 50]$ и не затрагивает неизвестную внутреннюю область $s < e(t) - W + 1$) и аналогично к левому фронту $s = -t$ с шагом $(\Delta t, \Delta s) = (40,-40)$, получаем первую крупную бесконечную теорему.

\begin{theorem}[Две приграничные полосы]\label{thm:dual_strips}
Для каждого целого $t \ge 300$ выполнено
\begin{equation}\label{eq:dual_strips}
[-L(t), -L(t) + 105] \cap \mathbb{Z} \;\cup\; [t - 7, t] \cap \mathbb{Z} \;\subseteq\; D_t.
\end{equation}
\end{theorem}

Левая полоса ширины $106$ клеток сертифицируется $20$ фазами фонов $(0001111001)$, $(1011)$ и $(001011)$ с глубиной тайла $167$ клеток на $30$ базовых тактах $t = 300, \dots, 329$. Правая полоса ширины $8$ клеток сертифицируется шестью фазами с шагом $(\Delta t, \Delta s) = (40,-40)$ на базовых тактах $t = 300, \dots, 339$. Происхождение промежуточных состояний удостоверено по схеме \emph{proof-carrying computation} на битовых векторах \texttt{BitVec 1320} с проверкой каждого перехода $F_{1320}^{h_i}(z_i) = z_{i+1}$ ($h_i \in \{1, 10\}$) в ядре Lean 4.35.0.

\subsection{Баллистические лучи и границы локальных сертификатов}

Второй механизм неограниченного продолжения во времени основан на сохранении различия вдоль дискретных мировых линий.

\begin{lemma}[Сохранение правого полулуча]\label{lem:ballistic_ray}
Если для траектории $A$ в момент $T$ выполнено $\forall s \ge s,\; A_{T+3}(s + 2) = A_T(s)$, то $\forall q \ge 0,\; \forall s \ge s,\; A_{T+3q}(s + 2q) = A_T(s)$. В частности, если это условие выполнено одновременно для $B$ и $\widetilde B$ и $B_T(s) \ne \widetilde B_T(s)$, то $B_{T+3q}(s + 2q) \ne \widetilde B_{T+3q}(s + 2q)$ для всех $q \ge 0$.
\end{lemma}

Для двух арифметических серий якорей на фонах $(0000101111)$ и $(001011)$ при $0 \le n \le 32$:
\begin{equation}
T_A(n) = 594 + 882n,\quad s_A(n) = 256 + 378n; \qquad
T_B(n) = 884 + 882n,\quad s_B(n) = 380 + 378n,
\end{equation}
вычислительно сертифицированы $66$ начальных точек, порождающих бесконечные баллистические лучи существенности на постоянных вдоль каждого луча глубинах $d_A(n) = 140 + 210n$ и $d_B(n) = 210 + 210n$ от фронта $e(t) = L(t)$.

В последующих модулях \texttt{StaircaseA} и \texttt{StaircaseB} формулы обеих серий доказаны для произвольных $n,q\ge0$. Для этого три локально согласованные строки склеиваются в точную эволюцию. На передних периодических средах серий A и B обнаружено $45\,255$ покрытых классов из $66\,150$; с третьей серией --- $53\,550$. Число классов относится к таблице остатков $315\times210$. Применение таблицы в физической точке требует также проверки границ соответствующего клина.

Важную методологическую роль сыграли строгие отрицательные результаты, уточнившие границы применимости простых гипотез:
\begin{itemize}
    \item \textbf{Опровержение локально-тайловой макроиндукции по $n \mapsto n+1$:} при переходе на макрошаг $(\Delta T, \Delta s) = (882, +378)$ универсальная область зависимости составляет $\mathrm{Dep}_{882}(s) \subseteq [s - 588, s + 882]$ (прообраз $[s - 210, s + 1260]$). Прямыми контрпримерами в \texttt{check\_macro\_boundary.py} показано, что фиксация локальных тайлов радиуса $64$ (для $T=594$) или $186$ (для $T=884$) недостаточна: изменение клеток вне тайла меняет целевой бит через $882$ такта.
    \item \textbf{Неполнота наивных угловых секторов:} для фона $(0001111001)$ наблюдаемый до $t = 2000$ левый сектор со скоростями $[-0.66, -0.64]$ обнаруживает первый пропуск на более длинном горизонте.
    \item \textbf{Разделение периодичности и полноты покрытия:} глубокая сдвиговая рекуррентность ($261$ клетка для $(0000001001)$) не гарантирует сплошного покрытия без объединения нескольких фаз по всем классам вычетов.
\end{itemize}

\subsection{Конструктивное замыкание бесконечного левого клина: теорема \texttt{wedge\_all}}

Чтобы преодолеть несоизмеримость скоростей между левым краем светового конуса (скорость $-1$), фоновым кристаллом периода $50$ и фронтом распространения возмущения, вводится дискретная параметризация чётных тактов $n \in \mathbb{N}$:
\begin{equation}
\mathrm{tm}(n) = 2450 + 2n, \qquad a(n) = -350 - n.
\end{equation}
Здесь правая граница клина $a(n)$ смещается влево со скоростью $-1/2$ клетки за такт ($1$ клетка за $2$ такта). Для каждой из шести фаз $p \in \{\mathrm{p0}, \mathrm{p1}, \mathrm{p2}, \mathrm{p3}, \mathrm{p7}, \mathrm{p9}\}$ пишем $\widetilde B_p(t,s)=\widetilde B^p_t(s)$. Сопряжённый инвариант на слое $n$ объединяет совпадение с фоновым кристаллом $\mathrm{Cm}_p$ на интервале слева от $a(n)$ и принадлежность $16$-битного граничного буфера узлу конечного графа переходов:
\begin{align}
\mathrm{Medium}(p,n)\ \iff\ {}&
\forall s\in[-\mathrm{tm}(n)+1300,a(n))\cap\mathbb{Z},\notag\\
&\widetilde B_p(\mathrm{tm}(n),s)=\mathrm{Cm}_p(\mathrm{tm}(n),s), \\
\mathrm{Window}(p,n,i)\ \iff\ {}&
\forall k\in\{0,\dots,15\},\notag\\
&\widetilde B_p(\mathrm{tm}(n),a(n)+k)=\mathrm{testBit}(\mathrm{word}(p,n,i),k).
\end{align}

\begin{lemma}[16-битный детерминированный трансдьюсер и склейка со швом]\label{lem:transducer}
Для каждой фазы $p$ таблица \texttt{Crystal50Packed} задаёт $50$-периодический граф $36$-битных записей (младшие $16$ бит кодируют окно $\mathrm{word}(p, n, i)$, биты $16\dots 25$ и $26\dots 35$ — переходы при правом бите $r \in \{0, 1\}$).
\begin{enumerate}[label=(\alph*)]
    \item На шаге $n$ $4$-битный левый контекст $\mathrm{ctx}(p, n \bmod 50)$ в точках $a(n) - 4 + m$ ($m \in \{0, 1, 2, 3\}$) имеет глубину $t + s = 2096 + n + m$ и принадлежит уже доказанному интервалу $\mathrm{Medium}(p, n)$.
    \item Двукратное побитовое сжатие $\mathrm{shrink}(\mathrm{shrink}(\mathrm{input}(p, n, i, r)))$ вычисляет $17$ бит конфигурации в момент $\mathrm{tm}(n+1)$ на $[a(n)-2, a(n)+14]$: младший бит ($k=0$) даёт $\mathrm{Cm}_p(\mathrm{tm}(n+1), a(n) - 2)$, а старшие $16$ бит ($k=1\dots 16$) совпадают с узлом $z$ следующего слоя графа в позиции $a(n+1) = a(n) - 1$.
    \item Левый край полосы $-\mathrm{tm}(n+1) + 1300 \le s < -\mathrm{tm}(n) + 1302$ имеет инвариантную глубину $k = s + \mathrm{tm}(n+1) \in \{1300,1301,1302,1303\}$ и закрывается теоремой \texttt{seam} (\texttt{Crystal50Seam}) благодаря инвариантности $t + s$ при сдвиге $(t, s) \mapsto (t + 50, s - 50)$.
\end{enumerate}
\end{lemma}

\begin{theorem}[\texttt{Rails.Crystal50Wedge.wedge\_all}]\label{thm:wedge_all}
Для любого натурального $t \ge 2485$ и любого целого $s$, удовлетворяющего неравенствам
\begin{equation}\label{eq:wedge_region}
-t + 1302 \le s \le 59 - \lfloor t/2 \rfloor,
\end{equation}
выполняется предикат \texttt{EssentialAt}\,$t\,(-s)$. Верхней границы по времени нет.
\end{theorem}

\begin{proof}[Схема доказательства]
Положим $n = \lfloor (t - 2450) / 2 \rfloor$. При $t \ge 2485$ имеем $n \ge 17$, а правая граница $59 - \lfloor t/2 \rfloor = -1166 - n = a(n) - 816$. Для чётных тактов $t = \mathrm{tm}(n)$ равенство $\widetilde B_p(t, s) = \mathrm{Cm}_p(t, s)$ следует непосредственно из инварианта $\mathrm{Medium}(p, n)$ для всех шести фаз $p$. Для нечётных тактов $t = \mathrm{tm}(n) + 1$ одношаговый конус зависимости $[s - 1, s + 1]$ целиком лежит внутри $[-\mathrm{tm}(n) + 1300, a(n))$, и совпадение переносится леммой локальности \texttt{carry}. Применение моста \texttt{essential\_of\_cm} завершает доказательство.
\end{proof}

\subsection{Полное покрытие на бесконечной прямой}\label{sec:five-bands}
В координате входа $j=-s$ положим $b=\lfloor t/2\rfloor$. При $t\ge2485$ последовательно применяются следующие пять результатов. Обе границы в каждой строке включены.
\begin{center}
\begin{tabular}{lll}
\toprule
Начало & Конец & Основание \\
\midrule
$-L(t)$ & $-L(t)+105$ & Полулучевой сертификат ширины 106 \\
$-L(t)+106$ & $b-60$ & \texttt{EtherWedge.half\_sector} \\
$b-59$ & $t-1302$ & \texttt{Crystal50Wedge.wedge\_all} \\
$t-1301$ & $t-20$ & \texttt{LeftCrystal50.left\_strip} \\
$t-19$ & $t$ & \texttt{right\_boundary\_width20} \\
\bottomrule
\end{tabular}
\end{center}
Три фазы эфира периода 14 дают вторую строку уже при $t\ge210$. Последняя строка усиливает раннюю правую полосу ширины восемь. При $t=2485$ третья строка состоит из единственной клетки $j=1183$. Начало каждой следующей строки на единицу больше конца предыдущей. Лемма \texttt{Rails.late\_all} объединяет пять доказательств.

\paragraph{Конечный сертификат.}
Для $0\le t<2485$ берём все циклические фазы семи слов. Их длины и целые кодировки (младший бит соответствует первой клетке) равны
\[
(8,244),\ (10,632),\ (10,976),\ (10,576),\ (4,13),\ (6,52),\ (6,48).
\]
Фоны представлены векторами длины 5520; однобитный дефект внесён в позицию 2760. Периоды всех семи слов делят 5520. Для каждой строки времени строим побитовое объединение различий всех пар и проверяем, что оно содержит маску физического интервала $[-t,L(t)]$. Это $\sum_{t=0}^{2484}(2t+1-\lfloor t/3\rfloor)=5\,147\,263$ пары. При $s$ внутри интервала выполнено $-5520<s-t$ и $s+t<5520$. Поэтому лемма \texttt{local\_obs} переносит результат на бесконечную прямую с единственным переворотом.

Начальная строка вычисляется из семи слов. Генератор предлагает промежуточные снимки с шагом 25. Lean проверяет маски каждого блока и равенство его конечного состояния начальному снимку следующего. Последний блок содержит десять строк. Индукция по длине блока связывает бит маски с конкретной парой траекторий. Таким образом, лемма \texttt{Rails.Finite.finite\_all} доказывает существенность для каждого входа конечной области.

\begin{proof}[Завершение доказательства теоремы~\ref{thm:gap-free}]
Включение $D_t\subseteq I_t$ следует из теоремы~\ref{thm:upper_bound}. Для обратного включения применяем конечный сертификат при $t<2485$ и пять полос при $t\ge2485$. Каждый свидетель различия в физической координате $s$ переносится сдвигом в свидетельство существенности входа $-s$.
\end{proof}

\paragraph{Следствие для доступа к данным.}
Любое фиксированное множество исходных координат, достаточное для точного вычисления $G_t$ на всех конфигурациях, содержит $I_t$. Достаточность самого $I_t$ даёт локальность вместе с верхней границей. Минимальное число координат такого наблюдения равно $2t+1-\lfloor t/3\rfloor$. Экзистенциальный свидетель может зависеть и от $t$, и от $j$. Для выводов о стоимости адаптивного алгоритма, предобработке или объёме памяти требуется отдельно задать модель вычислений.

Имена модулей, область применимости каждого машинного сертификата и связь с теоремами о наблюдателе приведены в приложении~\ref{app:lean}.
